\documentclass[10pt]{article}

\usepackage[margin=0.85in]{geometry}
\usepackage[T1]{fontenc}
\usepackage[utf8]{inputenc}
\usepackage{amsmath,amssymb,booktabs,array,microtype}
\usepackage{xcolor}
\usepackage{hyperref}
\usepackage{url}

\hypersetup{
  colorlinks=true,
  linkcolor=blue!55!black,
  citecolor=blue!55!black,
  urlcolor=blue!55!black
}

\newcommand{\tikitorch}{\textsc{TikiTorch}}
\newcommand{\todo}[1]{\textcolor{red!70!black}{\textbf{TODO:} #1}}
\newcommand{\systemname}{Multiplier-Free Learning on FPGA}

\title{\systemname:\\
Closing the Forward--Backward--Update Loop on Real Silicon}
\author{Anonymous authors\\
\small Author list and affiliations to be finalized}
\date{Internal draft --- August 2026}

\begin{document}
\maketitle

\begin{center}
\fcolorbox{red!55!black}{red!4}{%
\begin{minipage}{0.93\linewidth}
\textbf{Draft status.}
This document is an evidence-bounded paper scaffold, not a submission-ready
manuscript.  Reported numbers are tied to committed artifacts.  Missing
iso-model CUDA, disjoint-stream, timing-report, and full-card-power experiments
are marked explicitly and must not be inferred from the present results.
\end{minipage}}
\end{center}

\begin{abstract}
Accelerators for low-precision and multiplication-less neural networks usually
optimize inference or selected linear layers, while training continues to rely
on conventional products elsewhere in the forward pass, backward pass, loss,
or optimizer.  We study a stricter systems question: can a useful neural
network execute its complete training step on field-programmable gate array
(FPGA) fabric without data-dependent multiplication?  We co-design hard-routed
tropical and dyadic operators with sparse reverse routing, lookup-table or
subtractive losses, and integer sign-based updates.  A three-level closure
auditor checks model structure, an executed operator trace, and the full
training recipe before synthesis.

On an AWS F2 FPGA, a four-layer, width-512 morphological state-space language
model executes forward, loss, backward, and SignSGD in one fused invocation.
The 30-step silicon result is bit-for-bit identical to the C-simulation golden
result, and a 2,000-step invocation processes 256,000 token positions at 612
token/s.  This run repeats one resident 128-token window and therefore
establishes execution and persistent weight evolution, not language-model
generalization.  Separately, a two-layer dyadic network trained entirely on
fabric reaches 74.75\% test accuracy on 16-by-16 MNIST; an independent software
path evaluates the exported fabric-trained weights at 74.13\%.  These results
show that a complete multiplier-free learning loop can execute exactly on real
silicon and can produce weights that generalize on a small task.  They do not
yet establish an energy or throughput advantage over a matched GPU system;
those comparisons are left as required submission gates.
\end{abstract}

\section{Introduction}

Multiplication is deeply embedded in contemporary training systems: matrix
products dominate linear layers, activation derivatives multiply cotangents,
weight gradients form outer products, losses use nonlinear normalization, and
adaptive optimizers rescale updates.  Prior work has replaced some of these
operations with additions, shifts, or low-precision logic
\cite{deepShift,addernet,shiftaddnet,liu2023multfree}.  An implementation can
nevertheless be ``multiplication-less'' in one layer while retaining products
elsewhere in the training loop.  Conversely, source code that appears closed
can synthesize hidden multipliers for address arithmetic or an ostensibly
constant expression.

This paper treats multiplier freedom as a property of a complete executable
training recipe and its synthesized implementation.  The recipe combines
max/min-plus and dyadic operators, hard routing, a sparse reverse pass, a closed
loss, and a sign-based optimizer.  The software implementation lives in
\tikitorch, a PyTorch-shaped tropical-learning library; the FPGA implementation
is a fused high-level-synthesis (HLS) kernel on an AWS F2 device.

The paper is organized around three deliberately separate questions:

\begin{enumerate}
  \item \textbf{Closure:} Does the forward--loss--backward--update recipe avoid
  products of two run-time data values, and does synthesis agree?
  \item \textbf{Exact execution and scale:} Does a nontrivial, width-512
  sequence model execute the fused recipe on silicon with the same state
  transition as its reference implementation?
  \item \textbf{Learning:} Can fabric-only updates produce weights that perform
  above chance on held-out examples?
\end{enumerate}

Our current evidence answers these questions with complementary artifacts:

\begin{itemize}
  \item A machine-auditable closure interface distinguishes a static structural
  check, an executed operator trace, and a full-training check.  It also
  documents where each check is necessary but not sufficient.
  \item A four-layer width-512 MorphSSM character model runs an entire fused
  learning loop on F2.  Thirty updates are bit-exact against C simulation; a
  single 2,000-update invocation sustains 612 token/s over 256,000 processed
  token positions.
  \item A two-layer dyadic classifier is trained from scratch on the same class
  of multiplier-free fabric operations and reaches 74.75\% on held-out
  downsampled MNIST, with a 74.13\% independent evaluation of the exported
  weights.
\end{itemize}

We do \emph{not} claim state-of-the-art accuracy, a new differentiation rule,
weight-transport-free learning, or a GPU efficiency win.  The experiments are
designed to establish an exact end-to-end feasibility result and to identify
the measurements required for a fair systems comparison.

\section{Scope and terminology}

\subsection{Multiplier-free is not synonymous with zero DSP blocks}

We call an executed learning loop \emph{multiplier-free} when no model-data
path computes a product of two values that are both run-time dependent.  The
allowed primitive set is
\[
  \{\text{integer add/subtract},\; \text{compare},\; \text{mux},\;
    \text{absolute value},\; \text{bit shift},\; \text{table lookup}\}.
\]
Multiplication by a signed power of two is a shift and sign selection.  A
binary or ternary gate is a mux.  We report HLS and routed-resource evidence in
addition to algebraic closure because source-level index expressions may be
implemented differently by a compiler.

``Zero DSP'' is a stronger physical observation, not the definition.  The
MorphSSM builds report zero DSP blocks.  The MNIST build binds one DSP block as
an adder but reports zero DSP multipliers; its model-data path is still
multiplier-free.  The paper therefore avoids using ``zero DSP'' as a synonym
for all results.

\subsection{Three evidence levels}

\begin{table}[t]
\centering
\caption{Closure evidence used in this work.  No single software check is
treated as a substitute for synthesis and silicon validation.}
\label{tab:closure-levels}
\small
\begin{tabular}{p{0.16\linewidth}p{0.35\linewidth}p{0.39\linewidth}}
\toprule
Level & What is checked & Principal limitation \\
\midrule
Structural & Registered modules are scanned for dense linear, convolution,
normalization, attention, gated activation, and related multiply-bearing
modules. & Raw arithmetic inside a custom forward method is invisible. \\
Executed graph & A representative forward pass is intercepted at the ATen
operator level; dense products, general elementwise products/divisions, and
transcendentals are recorded. & Coverage is input dependent.  A host dyadic
layer may call BLAS even when its exported hardware schedule is shift--add. \\
Full training & Forward, loss, backward, and optimizer step are exercised and
their operations classified. & The software trace does not prove how HLS maps
integer expressions or memory addressing. \\
Physical & HLS reports, routed design, known-answer tests, and silicon/reference
state comparisons. & Applies only to the compiled kernel, shape, flags, and
test cases under report. \\
\bottomrule
\end{tabular}
\end{table}

\section{Related work}

Tropical algebra replaces ordinary addition and multiplication with max/min
and addition.  It provides a mathematical language for piecewise-linear neural
maps and hard routing \cite{zhang2018tropical}.  This work uses that algebra as
a hardware co-design constraint rather than making a general claim that all
neural-network computation is tropical.

Binary and quantized FPGA systems such as FINN demonstrate highly efficient
inference dataflows \cite{umuroglu2017finn}.  DeepShift represents weights using
shifts and signs \cite{deepShift}; AdderNet replaces convolutional products by
$\ell_1$-distance accumulation while retaining other conventional training
components \cite{addernet}; and ShiftAddNet combines shift and addition layers
with FPGA-backed measurements \cite{shiftaddnet}.  Ultra-low-precision
multiplier-free training replaces full-precision linear-layer multiplication
with low-precision operations \cite{liu2023multfree}.  These works are the
closest systems context, and they make a priority claim inappropriate without
a more systematic operation-by-operation and silicon-path comparison.

Our narrower distinction is empirical: we archive and compare the state
transition of a fused, complete training invocation---forward, loss, backward,
and update---running on an FPGA.  The update uses the sign of accumulated
gradients, related to signSGD \cite{bernstein2018signsgd}, but our paper makes no
new convergence claim for sign optimization.

\todo{Complete a systematic related-work matrix before submission.  For every
closest system, record which of forward, loss, backward, optimizer, synthesis,
routing, on-silicon execution, and held-out learning are demonstrated.  Remove
or retain novelty language only after this audit.}

\section{Closed learning recipe}

\subsection{Hard tropical and dyadic operators}

A max-plus projection computes
\begin{equation}
  y_j = \max_i \left(x_i + w_{ji}\right), \qquad
  r_j = \operatorname*{arg\,max}_i \left(x_i + w_{ji}\right).
\end{equation}
Its forward implementation uses additions and comparisons.  Under a fixed,
deterministically tie-broken route, reverse propagation scatters a cotangent
only to the selected input and parameter.  The hardware stores route indices
needed by its reverse sweep; the theory and retained-state analysis of such
route traces are treated in a separate paper.

Dyadic linear maps constrain each nonzero weight to
\begin{equation}
  q(w) = s\,2^e, \qquad s\in\{-1,+1\},\quad e\in\mathbb{Z}.
\end{equation}
Multiplication by $q(w)$ is implemented as a sign selection and bit shift.
The MNIST kernel additionally quantizes activation and error factors so that
the gradient accumulation is an exponent addition followed by a shift.  This
is a hardware learning rule; eager PyTorch autograd through a dyadic layer is
not automatically multiplier-free.

\subsection{Loss and update}

The language-model kernel supports a balanced lookup-table cross-entropy
gradient and a two-hot perceptron gradient.  The reported D512 artifact uses
the LUT loss.  The table is staged with the model buffer because an earlier
silicon build exposed that a source-level constant-ROM initialization could be
lost during synthesis.  This failure is one reason the evidence chain includes
on-silicon known-answer comparisons.

The optimizer accumulates integer gradient information and applies a sign
update.  A power-of-two learning rate is a shift.  The MNIST kernel uses a
capped leaky sign counter as integer momentum.  Neither path implements Adam,
and no claim is made that the closed optimizer matches Adam at equal quality or
training budget.

\section{System design}

\subsection{Software closure and export}

The \tikitorch{} audit surface exposes structural, run-time, and training-level
checks.  It treats a clean structural scan as necessary rather than sufficient;
custom source arithmetic and the dyadic eager/hardware distinction are called
out in the API documentation.  The HLS kernels are checked separately through
C simulation, synthesis reports, routed builds, and silicon runs.

The current FPGA backend uses fused whole-training-step monoliths matched to a
compiled model shape.  It is not yet arbitrary per-operator FPGA dispatch or a
general compiler from every \tikitorch{} model.  The host stages data, invokes
the kernel, and reads results, while the forward, loss-gradient, reverse sweep,
and optimizer execute in fabric.

\subsection{MorphSSM language-model kernel}

The scaled language model has a 128-symbol vocabulary, sequence length 128,
four alternating dilation/erosion morphological layers, seven causal taps, and
channel dimension $D=512$.  Each layer performs a batch-shift normalization, a
depthwise morphological convolution, a max-plus projection, an additive bias,
and a residual update.  A max-plus head emits character logits.

The fused training mode executes the following sequence for every optimizer
step:
\begin{enumerate}
  \item clear the rolling activation stream;
  \item sweep positions $0\ldots127$ through forward and loss;
  \item reuse the stream as a reverse accumulator;
  \item sweep positions $127\ldots0$ through the sparse backward;
  \item apply the resident SignSGD update.
\end{enumerate}
Weights and gradient accumulators are resident on chip during a fused
invocation and are written back for observation.  At $D=512$, the design uses
URAM for the large resident arrays and two 4-MiB staging banks.

\subsection{Dyadic MNIST kernel}

The learning-quality experiment uses a two-layer dyadic MLP on 16-by-16
downsampled MNIST.  One kick processes a minibatch of 32 through shift--add
forward propagation, a subtractive top error, a reverse cascade using the
forward weights through transposed access, dyadic gradient accumulation, and
integer-momentum SignSGD.  Weights and momentum persist in on-card DDR across
kicks; the host supplies pixels and one-hot targets and periodically reads
weights for evaluation.

The transposed read is not a separately materialized weight buffer, but it
still reuses the forward weights.  We therefore make no biological
weight-transport-free or feedback-alignment claim.  For the matched routes,
the cascade computes the corresponding reverse-mode gradient rather than a
new gradient rule.

\section{Experimental evidence}

\subsection{Platform and protocol}

All silicon results reported here use an AWS EC2 F2 instance with a Xilinx
Virtex UltraScale+ VU47P FPGA.  Exact artifact identifiers and immutable source
commits are listed in Appendix~\ref{sec:artifact-ledger}.  ``Bit-exact'' means
that the archived observable output record---including predictions and weight
CRC where applicable---matches the C-simulation golden record exactly.  It
does not mean that an unobserved internal signal has been exhaustively compared.

\todo{Archive the final paper AFIs, HLS reports, post-route timing and
utilization reports, host binaries, compiler versions, build flags, and
machine-readable run manifests under one immutable release.  The D512 AFI
executed successfully, but its final routed report must be copied into the
public artifact bundle before submission.}

\subsection{Exact D512 execution}

Table~\ref{tab:d512} reports the D512 fused invocations.  The 30-step silicon
record and independently executed C-simulation record both contain resident
accuracy 0.421875, weight CRC 10673268721020562856, and the same first 16
predictions.  The archive therefore marks the run bit-exact.  Longer silicon
invocations sustain 612 token/s and produce distinct weight CRCs and
predictions, confirming that the update remains live.

\begin{table}[t]
\centering
\caption{Width-512, four-layer MorphSSM fused execution on F2.  Each optimizer
step processes the same resident 128-token window in the current harness.}
\label{tab:d512}
\small
\begin{tabular}{rrrrl}
\toprule
Updates & Token positions & Wall time (s) & Throughput (token/s) & Evidence \\
\midrule
30   & 3,840   & 6.2879   & 611 & bit-exact to C simulation \\
200  & 25,600  & 41.8356  & 612 & silicon; state changes \\
2,000& 256,000 & 418.3963 & 612 & silicon; one fused invocation \\
\bottomrule
\end{tabular}
\end{table}

This result is an exact execution and scaling demonstration, not a language
quality result.  Repeating one window can change its resident-window accuracy
without establishing held-out prediction or adaptation to a stream.

\subsection{Preliminary D128 replication ladder}

The committed D128 artifact reports 14,236 token/s for one fused MorphSSM
kernel, 26,522 token/s for two replicas, and 53,001 token/s for four replicas.
Every instance matched its reference record, and the model-data path synthesized
with zero DSP blocks.  The two- and four-replica measurements show approximately
linear scaling within that build family.

These values are preliminary context rather than a final comparative result.
The one- and four-kernel routed reports have small negative WNS values
($-0.020$ ns and $-0.173$ ns), the current comparison uses a different GPU
model and optimizer, and reported FPGA power is the Vccint core rail rather
than full-card power.  We therefore do not use this ladder to claim timing
robustness, GPU parity, or an energy advantage.

\begin{table}[t]
\centering
\caption{Preliminary D128 replica measurements.  The table reports execution,
not an iso-model or iso-quality accelerator comparison.}
\label{tab:d128}
\small
\begin{tabular}{rrrrl}
\toprule
Replicas & Aggregate token/s & Scaling & DSP & Qualification \\
\midrule
1 & 14,236 & 1.00$\times$ & 0 & measured, bit-exact; WNS $-0.020$ ns \\
2 & 26,522 & 2.00$\times$ reported & 0 & measured, both replicas exact \\
4 & 53,001 & 4.00$\times$ reported & 0 & measured, all exact; WNS $-0.173$ ns \\
\bottomrule
\end{tabular}
\end{table}

\subsection{Held-out learning on MNIST}

The MNIST experiment trains for five epochs over 60,000 examples, using 1,875
fabric kicks per epoch.  Test accuracy rises to 71.01\% after epoch one and
reaches 74.75\% after epoch five (Table~\ref{tab:mnist}).  After the first kick,
37,729 of 65,536 first-layer weights differ from initialization, providing a
direct state-change diagnostic.

\begin{table}[t]
\centering
\caption{Held-out accuracy of the fabric-trained dyadic classifier.  Epoch wall
time includes host weight readback and evaluation over 10,000 test examples and
is therefore not a training-throughput measurement.}
\label{tab:mnist}
\small
\begin{tabular}{rrr}
\toprule
Epoch & Test accuracy (\%) & Inclusive wall time (s) \\
\midrule
1 & 71.01 & 24.8 \\
2 & 74.30 & 24.6 \\
3 & 73.99 & 24.6 \\
4 & 74.79 & 24.5 \\
5 & 74.75 & 24.8 \\
\bottomrule
\end{tabular}
\end{table}

A separate Python path loads the exported integer weights and a freshly loaded
MNIST test set and obtains 74.13\%.  The 0.62 percentage-point difference is
consistent with the documented image-resize rounding difference between data
pipelines.  More importantly, both evaluations use held-out examples and show
that the fabric-produced weights encode a generalizing classifier.

This is a structural correctness demonstration, not an MNIST accuracy or
throughput result.  The images are downsampled, the model is small, predictions
are class-imbalanced, and the reported epoch time includes host evaluation.

\section{Required submission experiments}

The current artifact supports feasibility and exactness.  A publishable systems
claim requires the following locked comparisons.

\subsection{Matched performance comparison}

\todo{Implement an iso-model CUDA baseline for the same hard-routed model, loss,
numeric representation, batch/window shape, update rule, and number of token
positions.  Report end-to-end and kernel-only throughput, latency distribution,
peak memory, memory traffic where measurable, and compilation/setup costs.}

The existing nanoGPT comparison is useful context but is neither iso-model nor
iso-quality.  It must not be used as the primary speedup denominator.  At least
one conventional dense baseline should additionally be matched by validation
quality and parameter count, with its algorithmic differences made explicit.

\subsection{Disjoint-stream learning}

\todo{Replace the resident-window repetition harness with a sequence of disjoint
training windows.  Pre-register held-out evaluation, identical token order,
number of updates, initialization, and optimizer state.  Compare FPGA and
reference trajectories at periodic checkpoints and confirm exactness or bound
the first divergence.}

This experiment is the gate from ``the training state machine executes'' to
``the scaled model learns from an actual stream.''  It should include at least
three seeds and a no-update control.

\subsection{Timing and power}

\todo{Archive final post-route timing for every headline build and require
non-negative setup/hold slack with margin.  Re-run any D128 replica result whose
current WNS is negative before using it as a headline throughput result.}

\todo{Measure full-card idle and active power with external or authoritative
board-level instrumentation.  Separate shell/DDR/static power from dynamic
kernel power; use synchronized sampling windows and report uncertainty.  Until
then, retain only the statement that Vccint measurements are core-rail values
and make no board-to-board energy claim.}

\subsection{Accuracy and ablations}

Required ablations include loss choice, sign-momentum state, dyadic rounding,
resident versus streamed weights, route-cache storage, and host orchestration.
The MNIST experiment should be repeated with frozen preprocessing, seeds, and a
matched CPU/GPU implementation.  The language model should report held-out
bits per byte or bits per character on the exact on-fabric trajectory rather
than resident-window accuracy.

\section{Limitations}

\begin{itemize}
  \item The D512 experiment repeats one 128-token window.  It proves exact
  execution, persistent updates, and sustained fused invocation, not held-out
  language learning.
  \item The MNIST experiment is small and uses downsampled inputs.  Its value is
  that held-out accuracy is produced by fabric-only updates, not that 74.75\%
  is competitive.
  \item The FPGA kernels are shape-specialized monoliths.  This paper does not
  demonstrate a compiler for arbitrary \tikitorch{} programs.
  \item Closure is scoped to the compiled recipe.  Soft relaxations, ordinary
  normalization, dense heads, autograd through unrestricted dyadic weights,
  and conventional adaptive optimizers can reintroduce multiplication.
  \item The current D128 replica results have marginally negative timing slack.
  Successful executions do not replace a timing-closed release artifact.
  \item Core-rail measurements are insufficient for an energy comparison with
  GPU board power.
  \item The reverse path reuses forward weights.  We do not claim feedback
  alignment or biological freedom from weight transport.
  \item We have not completed a systematic prior-art audit sufficient to claim
  the first multiplier-free FPGA training loop.
\end{itemize}

\section{Conclusion}

The committed artifacts establish two facts.  First, a complete learning state
transition---forward, loss, backward, and update---can be expressed with
compare, add, shift, mux, and lookup operations and can execute bit-exactly on
real FPGA silicon at width 512.  Second, the same broad co-design principle can
train a small classifier whose exported weights generalize to held-out data.

The next result must be comparative rather than merely demonstrative: a
timing-closed release build, disjoint-stream learning, an iso-model CUDA arm,
and full-card power.  Until those gates close, the defensible contribution is
exact end-to-end multiplier-free learning on FPGA silicon, not superiority to
modern GPU training.

\appendix

\section{Artifact and claim ledger}
\label{sec:artifact-ledger}

This draft was prepared against repository head
\texttt{9d13474} (branch \texttt{pc-morphssm}).  The following claim-bearing
artifacts already have immutable commits.

\begin{table}[h]
\centering
\caption{Committed evidence used by this draft.}
\small
\begin{tabular}{p{0.23\linewidth}p{0.49\linewidth}p{0.17\linewidth}}
\toprule
Claim & Repository artifact & Commit \\
\midrule
D512 30-step bit-exact and 2,000-step execution &
\path{results/fpga_silicon/d512_verify/} & \texttt{ee3c775} \\
D128 single/replica measurements and qualifications &
\path{results/lm_shakespeare/fig_hallmark_lm.md} & \texttt{0b33218} \\
MNIST fabric-only training and independent evaluation &
\path{results/predictive_coding/ONFABRIC_TRAINING.md} and
\path{results/predictive_coding/onfabric_train/} & \texttt{115a494} \\
Closure-auditor implementation and scope &
\path{tikitorch/audit.py} & source at draft head \\
D512 fused kernel and buffer contract &
\path{tikitorch/hardware/hls/forge_kernels/forge_train_step_morph_lm.cpp} and
\path{forge_buffer_morph_lm.hpp} & source at draft head \\
\bottomrule
\end{tabular}
\end{table}

The AFGI identifiers currently recorded by the artifacts are
\texttt{agfi-003a8a2f2b1ba7df7} for the D512 MorphSSM result and
\texttt{agfi-0fdc38db9da6add74} for the MNIST training result.  AFIs are not
assumed to remain publicly loadable; the release must include build inputs and
reports sufficient to reconstruct them.

\section{Reproduction sketch}

The D512 launcher is
\path{scripts/run_morph_lm_d512_verify_f2.sh}; the MNIST recipe and driver are
listed in \path{results/predictive_coding/ONFABRIC_TRAINING.md}.  These paths
currently depend on an AWS account, FPGA tooling, and a compatible F2 image.

\todo{Replace this sketch with a versioned artifact README containing exact
container/toolchain versions, checksums, expected cost, expected wall time,
one-command C-simulation, synthesis-only, and on-board tracks.}

\bibliographystyle{plain}
\bibliography{refs}

\end{document}
